% =====================================================================
\chapter{Autoevaluación}
\label{ap:autoevaluacion}
% =====================================================================

\section{Puntuación por parte}

\begin{center}
\begin{tabular}{@{}llcc@{}}
\toprule
Parte & Contenido & Puntos & Nota obtenida \\
\midrule
A & Cuestiones ($20\times 1{,}5$) & 30 & \\
B1 & Retropropagación en un perceptrón de dos capas & 18 & \\
B2 & Optimización & 10 & \\
B3 & Aprendizaje supervisado & 12 & \\
B4 & Diferenciación automática & 10 & \\
B5 & Entrenamiento de redes profundas & 10 & \\
B6 & Estimación de parámetros, cómputo y memoria & 10 & \\
\midrule
 & \textbf{Total} & \textbf{100} & \\
\bottomrule
\end{tabular}
\end{center}

La nota sobre~10 es la puntuación total dividida por~10. La nota de la fase es
\[
0{,}4\cdot\text{examen}+0{,}3\cdot\text{proyecto}+0{,}2\cdot\text{ejercicios}+0{,}1\cdot\text{tests},
\]
y exige, además, al menos un 70\,\% en este examen.

\section{Puntuación por módulo}

La tabla reparte los 100 puntos entre los módulos. El problema~6 combina varios módulos y se contabiliza aparte.

\begin{center}
\small
\begin{tabular}{@{}lccccc@{}}
\toprule
Módulo & Cuestiones & Puntos A & Problemas & Puntos B & Total \\
\midrule
M07 Optimización & 1--5 & 7,5 & 2; 1(g) & 10 + 1 & 18,5 \\
M08 Aprendizaje supervisado & 6--9 & 6 & 3; 1(b) & 12 + 3 & 21 \\
M09 Redes neuronales & 11--13 & 4,5 & 1(a), (c)--(g) & 14 & 18,5 \\
M10 Diferenciación automática & 10, 14, 15 & 4,5 & 4 & 10 & 14,5 \\
M11 Redes profundas & 16--20 & 7,5 & 5 & 10 & 17,5 \\
Varios (M07, M09--M11) & --- & --- & 6 & 10 & 10 \\
\midrule
\textbf{Total} & & \textbf{30} & & \textbf{70} & \textbf{100} \\
\bottomrule
\end{tabular}
\end{center}

En el problema~1 se asignan a M08 los 3 puntos del apartado~(b) (gradiente de la entropía cruzada), a M07 un punto del apartado~(g) (suma frente a media) y a M09 los 14 restantes.

\section{Umbral de aprobado y puerta de salida de la fase}

\begin{itemize}[leftmargin=1.4em]
  \item \textbf{Examen de hito:} al menos 70 de 100 puntos.
  \item \textbf{Puerta de salida de la Fase~II:} además del examen, se requiere una puntuación mínima de $9/12$ en el test de cada módulo (M07 a M11) y al menos un 70\,\% en la autoevaluación del laboratorio interactivo de la fase.
  \item \textbf{Proyecto de fase:} debe superar al modelo de bigramas ($\PPL_{\text{val}}<12{,}42$) para obtener al menos la calificación de suficiente en el criterio de resultado.
\end{itemize}

\section{Qué hacer si no se alcanza el umbral}

Si un módulo obtiene menos del 70\,\% de sus puntos en la tabla anterior, conviene repasar las secciones indicadas, rehacer los ejercicios $\star\star\star$ de esos temas en los apuntes del módulo y repetir el examen al cabo de una semana sin consultar las soluciones.

\begin{center}
\small
\renewcommand{\arraystretch}{1.12}
\begin{tabularx}{\textwidth}{@{}>{\raggedright\arraybackslash}p{4.2cm}>{\raggedright\arraybackslash}X@{}}
\toprule
Si se falla\dots & Repasar \\
\midrule
Cuestiones 1--5, problema~2 & M07, §2.2--2.3 (convergencia y $\kappa$), §3.1 (minibatch), §4.1--4.2 (momentum), §5.3 (Adam), §6.2--6.3 (AdamW, calentamiento y coseno) \\
Cuestiones 6--9, problema~3 & M08, §2.2 y §2.5 (ecuaciones normales y máxima verosimilitud), §3.4 (logística), §4.3 y §4.5 (softmax y capa de salida), §5.3 (diagnóstico), §6.5 (validación cruzada) \\
Cuestiones 11--13 & M09, §2.2--2.4 (ReLU, GELU y SwiGLU), §4.3 (Xavier) \\
Problema~1 & M09, §5.2--5.5 (retropropagación y ejemplo completo), §5.6 (comprobación numérica); M08, §4.3 \\
Cuestiones 10, 14, 15; problema~4 & M10, §1.1 (trazas), §2.2--2.3 (números duales y coste), §3.1--3.2 (adjuntos), §4.4 (regla $6ND$), §5.1 (memoria), §6.2 (PyTorch) \\
Cuestión 16, problema~5(b) & M11, §3.1--3.2 (LayerNorm, RMSNorm e invariancias) \\
Cuestión 17, problema~5(c) & M11, §1.3 (saturación), §2.2 (residuales), §4.1--4.2 (dropout y recorte) \\
Cuestiones 18--19, problema~5(d) & M11, §5.1--5.3 (formatos, redondeo y precisión mixta) \\
Cuestión 20, problema~5(a) & M11, §6.1--6.4 (log-sum-exp y diagnóstico de \NaN) \\
Problema~6 & M09, §6.3 (modelo de Bengio); M10, §4.4; M07, §5.3; M08, §4.5 \\
Proyecto por debajo del bigrama & M09, §6.3--6.4; M10, §6.2; M07, §6.2--6.3; M11, §4.1 \\
\bottomrule
\end{tabularx}
\end{center}

\paragraph{Siguiente.} Fase~III, módulo~12: tokenización.
